\chapter{Diagnóstico: qué hacer si algo no anda}\label{cap:diag}

\section{Primero, mirar el estado}
Casi todo problema aparece en el texto de estado de cada placa:
\begin{lstlisting}[language=consola]
ros2 topic echo /base/estado --field data       # modo, seguridad, bloqueos, IMU, fallas, enlace
ros2 topic echo /brazo/estado --field data      # armado; por articulación: modo, cero, duty, falla, llego
ros2 topic echo /base/respuesta --field data    # respuesta a cada comando (por qué se rechazó)
ros2 topic echo /brazo/respuesta --field data
ros2 topic echo /mision/estado --field data     # estado y motivo de la misión
ros2 topic echo /brazo/ik_estado --field data   # por qué un punto es inalcanzable
\end{lstlisting}

\section{Tabla de síntomas}

\begin{longtable}{@{}p{4.6cm} p{5.0cm} p{5.9cm}@{}}
\toprule
\textbf{Síntoma} & \textbf{Causa probable} & \textbf{Qué hacer} \\
\midrule
\endhead
\multicolumn{3}{@{}l}{\textbf{Enlace y micro-ROS}} \\
\topic{/odom} o \topic{/brazo/joint\_states} no aparecen & Agente no corriendo, puerto equivocado o placa sin hora &
  Arrancar el agente con el puerto de \cod{/dev/serial/by-id/}; esperar unos segundos \\
La S3 no abre sesión & Quedó en modo descarga & \cod{esptool --before no\_reset --after watchdog-reset read-mac} \\
El brazo se reconecta cada $\sim$1~s & Baudios altos en el CH340 & Usar 115\,200 (el agente con \cod{-b 115200} si se cambió) \\
\cod{pings\_fallidos} o \cod{cortes} suben & USB flojo o la Pi sin energía & Revisar cables; \cod{dmesg}; \cod{vcgencmd get\_throttled} \\
No se puede flashear & El agente tiene el puerto & Detener el agente antes de \cod{pio run -t upload} \\
\midrule
\multicolumn{3}{@{}l}{\textbf{Base}} \\
\topic{/cmd\_vel} no mueve nada & Base desarmada & \cod{armar} por \topic{/base/comando} \\
Armada pero no se mueve & Seta, vuelco, IMU sin datos, sin eco al frente, obstáculo cerca & Leer \topic{/base/estado}; despejar la causa y \cod{r} \\
Avanza a los tirones y frena & \cod{/cmd\_vel} publicado con \cod{--once} o teleop sin \cod{repeat\_rate} & Publicar a 10~Hz (\cod{-r 10}; \cod{repeat\_rate:=10.0}) \\
Una rueda gira al revés & Signo del motor & \cod{MOTOR\_INVERTIDO} \\
rpm negativas yendo adelante & Signo del encoder & \cod{ENC\_INVERTIDO} \\
Avanza menos o más que lo pedido & Radio de rueda o cuentas por vuelta & Medir con cinta y corregir \cod{CFG\_DIF} \\
Gira de más o de menos & Trocha & Medir y corregir \cod{trocha} \\
Falla de rueda / atasco & Rueda trabada o encoder desconectado & Revisar; \cod{r} \\
Encoder cuenta solo, quieto & Entrada al aire (falta pull-up) o ruido & Revisar cableado y pull-ups \\
\cod{BATERIA\_BAJA} con la batería llena & Seta apretada (divisor en 0~V) o divisor mal & Normal con la seta; si no, medir el divisor \\
Bumper o seta disparan solos & Ruido (no hay capacitores) & 100~nF de ese pin a GND \\
\midrule
\multicolumn{3}{@{}l}{\textbf{Brazo}} \\
Rechaza movimientos & Desarmado o sin cero & \cod{armar}; \cod{z} en la pose de referencia o \cod{home <n>} \\
\cod{FueraDeLimites} & Se movió a mano o el cero está mal & \cod{r}; rehacer el cero \\
Atasco & Obstáculo, correa floja o encoder desconectado & Revisar; \cod{r} \\
Tiembla cerca del objetivo & PI o $v_\text{mín}$ mal ajustados & Bajar $K_p$; ajustar \cod{velMin} y tolerancia \\
J2 o J3 caen & No se autobloquean & Apoyar el brazo; nunca \cod{l} con el brazo levantado \\
\cod{inalcanzable} en \topic{/brazo/ik\_estado} & Punto fuera del espacio de trabajo o del rango de J1 & Leer el motivo; acercar la base \\
Servo da un salto al encender & GPIO~14 pulsa al arrancar & Conectar los 6~V después de que arranque la WROOM, o pasar la señal al GPIO~13 \\
\midrule
\multicolumn{3}{@{}l}{\textbf{Visión y misión}} \\
Sin \topic{/limon/posicion} con el limón a la vista & Limón en la zona ciega, fuera de 0,38--0,96~m o robot en movimiento & Correr el robot; frenar \\
Posiciones que saltan & Exposición sin sincronizar, robot moviéndose & Localizar quieto \\
Cámara no arranca & Cinta CSI & Intercambiar cintas con la Pi apagada \\
Colores rosados al sol & Cámaras NoIR & Pendiente: probar o sombrear \\
Misión en FALLA & Ver el motivo en \topic{/mision/estado} & Corregir y reiniciar el nodo \\
Misión se queda en ESPERAR & Limón inestable, brazo no quieto o fuera de alcance & Ver \topic{/mision/estado}; acercar la base \\
La simulación mueve el robot real & Falta aislar & \cod{ROS\_AUTOMATIC\_DISCOVERY\_RANGE=}\allowbreak\cod{LOCALHOST} \\
\midrule
\multicolumn{3}{@{}l}{\textbf{Energía}} \\
La Pi se reinicia o \cod{get\_throttled} $\neq$ \cod{0x0} & Power bank al límite o cable largo & Una sola salida, cable corto; bajar carga \\
Salta el fusible & Cortocircuito & \textbf{No} reemplazar sin encontrar el corto \\
\bottomrule
\end{longtable}

\section{Comandos útiles de la Pi}

\begin{lstlisting}[language=consola]
vcgencmd get_throttled                  # 0x0 = sin problemas de alimentación ni temperatura
vcgencmd pmic_read_adc EXT5V_V          # tensión de entrada (>= 4,9 V)
dmesg -w                                # desconexiones USB
ls -l /dev/serial/by-id/                # qué placa está en qué puerto
htop                                    # uso de CPU (YOLO + SGBM son lo más pesado)
\end{lstlisting}
